% =============================================================================
% 08b-cpu-scheduling-all-algorithms.tex — Special Lecture:
% CPU Scheduling — The Complete Toolkit
%
% Course: Operating Systems (PCC CS-502)
% A consolidation/revision deck after Lectures 8 and 9: every CPU-scheduling
% algorithm on the GATE/UG syllabus, one measurement vocabulary, worked
% examples on a single workload W, and a GATE-pattern homework set with
% full solutions.
%
% Anchor reading: Silberschatz, Galvin & Gagne, Operating System Concepts,
%   8th ed. — Ch. 5 §5.1–5.5 (pp. 183–206); Stallings, Operating Systems:
%   Internals and Design Principles, 6th ed. — Ch. 9 §9.1–9.2 (pp. 406–420);
%   Bovet & Cesati, Understanding the Linux Kernel, 3rd ed. — Ch. 7.
%
% Every Gantt chart and every CT/TAT/WT/RT number in this deck was produced
% by a discrete-event simulator and cross-checked by hand. Conventions:
%   * ties in SJF/priority: earlier arrival, then lower index;
%   * SRTF preempts only on a STRICTLY shorter remaining time;
%   * RR: an arrival at the instant a quantum expires joins the ready queue
%     BEFORE the preempted process;
%   * MLFQ: new arrivals enter Q0, a full quantum demotes one level, and a
%     higher queue preempts a lower one.
%
% Compile with: /compile-latex CS502/08b-cpu-scheduling-all-algorithms
% =============================================================================

\documentclass[aspectratio=169]{beamer}
\input{header}
\coursecode{CS502}

% -----------------------------------------------------------------------------
% Per-process shades, so a process stays visually identifiable across
% preemptions. These are only tracking colours, not semantic ones.
% -----------------------------------------------------------------------------
\colorlet{gP1}{light-bg}
\colorlet{gP2}{primary-gold!30}
\colorlet{gP3}{positive!18}
\colorlet{gP4}{hi-slate!22}
\colorlet{gP5}{neutral!18}
\colorlet{gFg}{light-bg}
\colorlet{gBg}{neutral!22}
\colorlet{gQ0}{light-bg}
\colorlet{gQ1}{primary-gold!30}
\colorlet{gQ2}{hi-slate!22}

% -----------------------------------------------------------------------------
% Gantt helper (same convention as Lecture 9):
%   #1 = cm per ms
%   #2 = slices  name/start/end/fillcolour
%   #3 = arrivals name/time   (pass {} for none)
% Bars at y=0..0.7, tick labels just below, arrival arrows at -0.75 and
% arrival labels at -0.95.
% -----------------------------------------------------------------------------
\newcommand{\gantt}[3]{%
  \begin{tikzpicture}[x=#1cm,y=1cm,baseline={(0,0.35)}]
    \foreach \n/\s/\e/\c in {#2} {
      \draw[draw=primary-blue,fill=\c] (\s,0) rectangle (\e,0.7);
      \node[font=\scriptsize] at ({(\s+\e)/2},0.35) {\n};
      \node[font=\tiny,anchor=north] at (\s,-0.02) {\s};
      \node[font=\tiny,anchor=north] at (\e,-0.02) {\e};
    }
    \foreach \n/\t in {#3} {
      \draw[-{Latex[length=1.2mm]},draw=primary-gold,thick] (\t,-0.75) -- (\t,-0.42);
      \node[font=\tiny,text=primary-gold] at (\t,-0.95) {\n};
    }
  \end{tikzpicture}%
}

% Compact per-process metrics table: #1 = body rows
\newcommand{\metrichead}{%
  \toprule
  \textbf{Proc.} & \textbf{AT} & $\boldsymbol{\tau}$ & \textbf{CT} &
  \textbf{TAT} & \textbf{WT} & \textbf{RT} \\
  \midrule
}

\title{CPU Scheduling: The Complete Toolkit}
\subtitle{Special Lecture --- Every Algorithm, Worked, with GATE-Level Practice}
\author[Dr. Auqib Hamid Lone]{%
  \textbf{Dr. Auqib Hamid Lone}\\
  \small Assistant Professor in Computer Science \& Engineering
}
\institute[GCET Ganderbal]{%
  \small\textbf{PCC CS-502: Operating Systems}\\[0.3em]
  Department of Computer Science \& Engineering\\
  Government College of Engineering \& Technology (GCET), Ganderbal
}
\date[\today]{\small\today}

\begin{document}

\frame{\titlepage}

% =============================================================================
% ACT 0 --- THE CONTRACT
% =============================================================================
\section{The Contract}

\begin{frame}{Today's Contract}
  \begin{block}{One workload, every algorithm}
    Five processes, one CPU, and the same yardsticks. We schedule the same
    work by FCFS, SJF, SRTF, Priority, Round Robin, and HRRN --- then meet
    the two queue-based algorithms, MLQ and MLFQ.
  \end{block}
  \begin{itemize}
    \item By the end you must be able to: draw the Gantt chart, fill the
      CT/TAT/WT/RT table, and defend the algorithm choice.
    \item The last third is a \key{GATE-pattern homework set} with full
      solutions --- attempt before turning the page.
  \end{itemize}
  \vspace{0.15cm}
  \begin{alertblock}{The one promise}
    No formula will be stated without being used on a number.
  \end{alertblock}
\end{frame}

\begin{frame}{The Running Workload $W$}
  \centering
  \renewcommand{\arraystretch}{1.25}
  \begin{tabular}{lcccc}
    \toprule
    \textbf{Process} & \textbf{Arrival (ms)} & \textbf{Burst $\tau$ (ms)} &
    \textbf{Priority} & \textbf{Character} \\
    \midrule
    \texttt{P1} & 0 & 7 & 3 & batch compile \\
    \texttt{P2} & 2 & 4 & 1 & interactive, important \\
    \texttt{P3} & 4 & 1 & 2 & tiny utility \\
    \texttt{P4} & 5 & 4 & 4 & background daemon \\
    \texttt{P5} & 9 & 3 & 2 & short, low-importance \\
    \bottomrule
  \end{tabular}
  \vspace{0.3cm}
  \begin{itemize}
    \item Total work $=7+4+1+4+3=19$ ms, so a work-conserving schedule with
      no idle time ends at $t=19$.
    \item \key{Lower priority number $=$ higher priority} (Silberschatz's
      convention). Ties go to the earlier arrival, then to the lower index.
  \end{itemize}
  \muted{Queue-based algorithms use their own small examples; the six
  single-queue policies all use $W$ so their numbers are comparable.}
\end{frame}

% =============================================================================
% ACT 1 --- THE MEASUREMENT VOCABULARY
% =============================================================================
\section{The Measurement Vocabulary}

\transitionslide{Before asking ``which algorithm is better?'', we must agree on what we are measuring.}

\begin{frame}{Vocabulary I: The Five Time Stamps}
  \centering
  \begin{tikzpicture}[x=0.55cm,y=1cm,every node/.style={font=\small}]
    % Timeline: waiting 0-3, first CPU 3-7, waiting again 7-9, CPU 9-12.
    \draw[draw=neutral!45,fill=neutral!12] (0,0) rectangle (3,0.7);
    \node[font=\scriptsize] at (1.5,0.35) {waiting};
    \draw[draw=primary-blue,fill=gP1] (3,0) rectangle (7,0.7);
    \node[font=\scriptsize] at (5,0.35) {CPU burst};
    \draw[draw=neutral!45,fill=neutral!12] (7,0) rectangle (9,0.7);
    \node[font=\scriptsize] at (8,0.35) {waiting};
    \draw[draw=primary-blue,fill=gP1] (9,0) rectangle (12,0.7);
    \node[font=\scriptsize] at (10.5,0.35) {CPU};
    \foreach \t in {0,3,7,9,12}
      \node[font=\tiny,anchor=north] at (\t,0) {\t};
    % AT / ST / CT markers
    \draw[-{Latex[length=1.3mm]},primary-gold,thick] (0,1.45) -- (0,0.85);
    \node[font=\tiny,text=primary-gold] at (0,1.65) {AT};
    \draw[-{Latex[length=1.3mm]},primary-gold,thick] (3,1.05) -- (3,0.85);
    \node[font=\tiny,text=primary-gold,anchor=south west] at (3,1.0) {ST};
    \draw[-{Latex[length=1.3mm]},primary-gold,thick] (12,1.45) -- (12,0.85);
    \node[font=\tiny,text=primary-gold] at (12,1.65) {CT};
    % RT brace
    \draw[decorate,decoration={brace,amplitude=4pt},neutral]
      (3,-0.35) -- (7,-0.35) node[midway,below=5pt,font=\tiny] {RT};
    % TAT brace
    \draw[decorate,decoration={brace,amplitude=5pt},primary-blue]
      (0,-0.95) -- (12,-0.95) node[midway,below=6pt,font=\tiny] {TAT};
  \end{tikzpicture}
  \vspace{0.45cm}
  \footnotesize
  \begin{tabular}{lll}
    \toprule
    \textbf{Time stamp} & \textbf{Symbol} & \textbf{Meaning} \\
    \midrule
    Arrival time & AT & enters the ready queue \\
    Burst time & BT or $\tau$ & CPU time this process needs (its next burst) \\
    Start / first dispatch & ST & first instant it gets the CPU \\
    Completion time & CT & instant its final burst ends \\
    Remaining time & --- & burst left; the key to preemptive SJF \\
    \bottomrule
  \end{tabular}
  \vspace{0.1cm}
  \par
  \muted{Waiting can be split: before ST and again after a preemption.}
\end{frame}

\begin{frame}{Vocabulary II: The Three Per-Process Metrics}
  \begin{block}{Definitions}
    For process $P_i$ with arrival $\mathrm{AT}_i$, burst $\tau_i$,
    first dispatch $\mathrm{ST}_i$, completion $\mathrm{CT}_i$:
    \[
      \boxed{\;\mathrm{TAT}_i=\mathrm{CT}_i-\mathrm{AT}_i\;}
      \qquad
      \boxed{\;\mathrm{WT}_i=\mathrm{TAT}_i-\tau_i\;}
      \qquad
      \boxed{\;\mathrm{RT}_i=\mathrm{ST}_i-\mathrm{AT}_i\;}
    \]
  \end{block}
  \begin{itemize}
    \item \key{Turnaround time} is the whole lifetime: waiting $+$ execution
      $+$ every re-wait caused by preemption.
    \item \key{Waiting time} is only the time spent ready but not running:
      $\mathrm{WT}=\mathrm{TAT}-\tau$. It is \emph{not} $\mathrm{CT}-\mathrm{AT}$.
    \item \key{Response time} stops at the first dispatch. Always
      $\mathrm{RT}\le\mathrm{WT}$; they are equal for a non-preemptive policy
      when each process has one burst.
    \item Average is $\frac1n\sum_i$ for each metric; the identity
      $\sum_i \mathrm{WT}_i=\sum_i \mathrm{TAT}_i-\sum_i \tau_i$ catches
      arithmetic slips.
  \end{itemize}
\end{frame}

\begin{frame}{Worked Example: Read Every Metric Off One Gantt Chart}
  \centering
  \footnotesize
  \texttt{P1}$(0,5)$, \texttt{P2}$(2,3)$, \texttt{P3}$(4,1)$ as
  (AT, $\tau$), scheduled FCFS.\\[0.15cm]
  \gantt{0.7}{P1/0/5/gP1,P2/5/8/gP2,P3/8/9/gP3}{P1/0,P2/2,P3/4}
  \\[0.15cm]
  {\footnotesize
  \begin{tabular}{lrrrrrr}
    \metrichead
    \texttt{P1} & 0 & 5 & 5 & 5 & 0 & 0 \\
    \texttt{P2} & 2 & 3 & 8 & 6 & 3 & 3 \\
    \texttt{P3} & 4 & 1 & 9 & 5 & 4 & 4 \\
    \midrule
    \multicolumn{4}{@{}l}{\textbf{Average}} & 5.33 & 2.33 & 2.33 \\
    \bottomrule
  \end{tabular}}
  \vspace{0.1cm}
  \begin{itemize}
    \item \texttt{P2}: $\mathrm{TAT}=8-2=6$, so $\mathrm{WT}=6-3=3$ ---
      not $8-2=6$.
    \item CPU utilization over $[0,9]$ is $9/9=100\%$; throughput is
      $3/9=0.33$ jobs/ms.
  \end{itemize}
\end{frame}

\begin{frame}{The Other Three Measurement Terms}
  \begin{columns}[T]
    \begin{column}{0.48\textwidth}
      \begin{block}{CPU utilization}
        Fraction of time the CPU does useful work:
        \[
          U=\frac{\text{busy}}{\text{busy}+\text{idle}}.
        \]
        Dispatcher time is not useful work.
      \end{block}
    \end{column}
    \begin{column}{0.48\textwidth}
      \begin{block}{Throughput}
        Completed jobs per unit time:
        \[
          X=\frac{\text{jobs completed}}{\text{interval}}.
        \]
        A system rate, not a per-process time.
      \end{block}
    \end{column}
  \end{columns}
  \vspace{0.1cm}
  \begin{itemize}
    \item \key{Dispatch latency} $s$: time to stop one process and start the
      next (context switch, mode switch, cache/TLB effects).
    \item On $W$ (19 ms of work), FCFS uses 5 dispatches. With $s=0.5$ ms:
      elapsed $=21.5$ ms, so
      $U=19/21.5\approx 88.4\%$ and $X=5/21.5\approx 0.233$ jobs/ms.
    \item More preemption can lower $U$ and $X$ even when it wins on TAT,
      WT, or RT.
  \end{itemize}
\end{frame}

\begin{frame}{When the Scheduler May Choose, and the Tie Rules We Use}
  \footnotesize
  A scheduling decision is needed when a process:
  \begin{enumerate}
    \item blocks (running $\to$ waiting);
    \item is preempted by the timer (running $\to$ ready);
    \item wakes (waiting $\to$ ready);
    \item terminates.
  \end{enumerate}
  \begin{columns}[T]
    \begin{column}{0.48\textwidth}
      \begin{block}{Non-preemptive}
        Decisions only in cases 1 and 4. Once running, a process keeps the
        CPU until it blocks or exits. Here $\mathrm{WT}=\mathrm{RT}$ for a
        one-burst process.
      \end{block}
    \end{column}
    \begin{column}{0.48\textwidth}
      \begin{alertblock}{Preemptive}
        Decisions also in cases 2 and 3. Necessary for interactive response,
        but every preemption pays dispatch latency and forces the kernel to
        protect shared data.
      \end{alertblock}
    \end{column}
  \end{columns}
  \vspace{0.15cm}
  \begin{itemize}
    \item \key{Tie rules used in every number here:} earlier arrival first,
      then lower index.
    \item SRTF preempts only for a \emph{strictly} shorter remaining time.
    \item RR: an arrival at the instant a quantum expires joins the queue
      \emph{before} the preempted process.
  \end{itemize}
\end{frame}

\begin{frame}{The Algorithm Family Tree}
  \centering
  \scriptsize
  \renewcommand{\arraystretch}{1.18}
  \setlength{\tabcolsep}{3pt}
  \begin{tabular}{p{1.25cm}p{2.95cm}p{1.65cm}p{2.4cm}p{2.95cm}}
    \toprule
    \textbf{Policy} & \textbf{Key idea} & \textbf{Preemptive?} &
    \textbf{What it favours} & \textbf{Main weakness} \\
    \midrule
    FCFS & serve in arrival order & no & simplicity, no starvation &
    convoy effect, long WT \\
    SJF & shortest next burst & no & minimum average WT (batch) &
    burst unknown; starves long jobs \\
    SRTF & shortest remaining time & yes & minimum average WT (preemptive) &
    more switches; starves long jobs \\
    LJF / LRTF & longest burst / remaining & both & almost nothing &
    worst averages; severe starvation \\
    Priority & highest importance first & both & policy goals &
    starvation without ageing \\
    RR & fixed time quantum & yes (timer) & response time &
    WT and TAT rise; switch overhead \\
    HRRN & highest response ratio & no & long waiters + short jobs &
    needs service times known \\
    MLQ & fixed process classes & both & class isolation &
    no migration; lower classes starve \\
    MLFQ & classes learned by behaviour & yes & responsiveness + short jobs &
    many parameters; can be gamed \\
    \bottomrule
  \end{tabular}
\end{frame}

\transitionslide{First, the four policies that order a single queue: arrival, burst, remaining time, and importance.}

% =============================================================================
% ACT 2 --- SINGLE-QUEUE POLICIES
% =============================================================================
\section{Single-Queue Policies}

\begin{frame}{FCFS: First-Come, First-Served}
  \begin{block}{Definition}
    \key{First-Come, First-Served} gives the CPU to the process that entered
    the ready queue first. The queue is FIFO: arrivals join the tail, the
    scheduler takes the head.
  \end{block}
  \begin{itemize}
    \item \key{Non-preemptive}: a process runs until it blocks or finishes.
    \item \key{Starvation-free}: everyone waits only for those already ahead.
    \item Average WT depends on arrival order, not just on the bursts ---
      the same jobs in a different order can differ by a large factor.
    \item Implementation is a linked list with head and tail pointers.
  \end{itemize}
  \muted{Silberschatz et al., Ch.~5 \S5.3.1, pp.~188--189; Stallings,
  Ch.~9 \S9.2, p.~410.}
\end{frame}

\begin{frame}{FCFS on $W$: The Numbers}
  \centering
  \footnotesize
  \gantt{0.6}{P1/0/7/gP1,P2/7/11/gP2,P3/11/12/gP3,P4/12/16/gP4,P5/16/19/gP5}%
    {P1/0,P2/2,P3/4,P4/5,P5/9}
  \\[0.1cm]
  {\footnotesize
  \begin{tabular}{lrrrrrr}
    \metrichead
    \texttt{P1} & 0 & 7 & 7 & 7 & 0 & 0 \\
    \texttt{P2} & 2 & 4 & 11 & 9 & 5 & 5 \\
    \texttt{P3} & 4 & 1 & 12 & 8 & 7 & 7 \\
    \texttt{P4} & 5 & 4 & 16 & 11 & 7 & 7 \\
    \texttt{P5} & 9 & 3 & 19 & 10 & 7 & 7 \\
    \midrule
    \multicolumn{4}{@{}l}{\textbf{Average}} & 9.0 & 5.2 & 5.2 \\
    \bottomrule
  \end{tabular}}
  \vspace{0.1cm}
  \begin{itemize}
    \item \texttt{P3} needs 1 ms and waits 7 ms; \texttt{P5} needs 3 ms and
      waits 7 ms. Arrival order, not job size, decides.
  \end{itemize}
\end{frame}

\begin{frame}{The Convoy Effect}
  \begin{itemize}
    \item One CPU-bound process takes a long burst; many short I/O-bound
      processes queue behind it and their devices sit idle.
    \item When the long job finally blocks, the short jobs flash through and
      block on I/O --- now the CPU sits idle.
    \item Everyone moves at the speed of the slowest, and utilization falls
      on both the CPU and the devices.
  \end{itemize}
  \vspace{0.1cm}
  \begin{exampleblock}{Two-minute proof}
    All arrive at 0; bursts $(20,1,1)$.\\
    FCFS order $20,1,1$: WT $=0,20,21$, average $=13.67$ ms.\\
    SJF order $1,1,20$: WT $=0,0,2$, average $=0.67$ ms.
  \end{exampleblock}
  \muted{Silberschatz et al., Ch.~5 \S5.3.1, pp.~188--189.}
\end{frame}

\begin{frame}{SJF: Shortest-Job-First (Non-preemptive)}
  \begin{block}{Definition}
    \key{Shortest-Job-First} gives the CPU to the ready process with the
    smallest \emph{next CPU burst}. Ties are broken FCFS.
  \end{block}
  \begin{itemize}
    \item The accurate name is \key{shortest-next-CPU-burst}: only the next
      burst matters, not the total job length.
    \item Non-preemptive here: once chosen, a process runs its whole burst.
    \item Needs a prediction of $\tau$ --- the scheduler cannot read the
      future (we return to this).
    \item Can starve a long job if short jobs keep arriving.
  \end{itemize}
  \muted{Silberschatz et al., Ch.~5 \S5.3.2, pp.~189--192; Stallings calls
  it SPN, Ch.~9 \S9.2.}
\end{frame}

\begin{frame}{SJF on $W$: The Numbers}
  \centering
  \footnotesize
  \gantt{0.6}{P1/0/7/gP1,P3/7/8/gP3,P2/8/12/gP2,P5/12/15/gP5,P4/15/19/gP4}%
    {P1/0,P2/2,P3/4,P4/5,P5/9}
  \\[0.1cm]
  {\footnotesize
  \begin{tabular}{lrrrrrr}
    \metrichead
    \texttt{P1} & 0 & 7 & 7 & 7 & 0 & 0 \\
    \texttt{P2} & 2 & 4 & 12 & 10 & 6 & 6 \\
    \texttt{P3} & 4 & 1 & 8 & 4 & 3 & 3 \\
    \texttt{P4} & 5 & 4 & 19 & 14 & 10 & 10 \\
    \texttt{P5} & 9 & 3 & 15 & 6 & 3 & 3 \\
    \midrule
    \multicolumn{4}{@{}l}{\textbf{Average}} & 8.2 & 4.4 & 4.4 \\
    \bottomrule
  \end{tabular}}
  \vspace{0.1cm}
  \begin{itemize}
    \item At $t=7$ the ready set is $\{4,1,4\}$: \texttt{P3} (1 ms) runs
      first; \texttt{P2} beats \texttt{P4} on the arrival-time tie-break.
    \item Average WT falls from 5.2 ms (FCFS) to \good{4.4 ms}.
  \end{itemize}
\end{frame}

\begin{frame}{Why SJF Is Optimal (When All Jobs Are Ready)}
  \begin{itemize}
    \item $n$ processes ready at time 0, run in order with bursts
      $\tau_{(1)},\tau_{(2)},\dots,\tau_{(n)}$.
    \item The $k$-th burst delays every one of the $n-k$ processes after it:
      \[
        \sum_i \mathrm{WT}_i=\sum_{k=1}^{n}(n-k)\,\tau_{(k)}.
      \]
    \item The first burst carries weight $n-1$, the largest. Put the
      smallest burst there, then repeat on the rest.
  \end{itemize}
  \begin{alertblock}{Result and its fine print}
    For a fixed batch of simultaneously available jobs, non-preemptive SJF
    minimises average waiting time. The proof is an exchange argument:
    swapping an adjacent long--short pair into short--long lowers the total.
    With arbitrary arrivals, SJF is a heuristic, not a theorem.
  \end{alertblock}
  \muted{Silberschatz et al., Ch.~5 \S5.3.2, p.~189.}
\end{frame}

\begin{frame}{The Catch: Predicting the Next Burst}
  \begin{itemize}
    \item The scheduler can only predict the next burst from past ones. The
      standard predictor is the \key{exponential average}:
  \end{itemize}
  \begin{block}{Exponential averaging}
    \[
      \hat{\tau}_{n+1}=\alpha\,\tau_n+(1-\alpha)\,\hat{\tau}_n,
      \qquad 0\le\alpha\le1.
    \]
    $\alpha=0$: history only. $\alpha=1$: only the last burst. Expanding,
    older bursts get weights $\alpha(1-\alpha)^j$.
  \end{block}
  \vspace{0.05cm}
  \centering
  \scriptsize
  \begin{tabular}{lcccccccc}
    \toprule
    $n$ & 0 & 1 & 2 & 3 & 4 & 5 & 6 & 7 \\
    \midrule
    measured $\tau_n$ & 6 & 4 & 6 & 4 & 13 & 13 & 13 & --- \\
    predicted $\hat{\tau}_n$ & 10 & 8 & 6 & 6 & 5 & 9 & 11 & 12 \\
    \bottomrule
  \end{tabular}
  \vspace{0.1cm}
  \muted{$\hat{\tau}_0=10$, $\alpha=1/2$. Example follows Silberschatz
  et al., Ch.~5 \S5.3.2.}
\end{frame}

\begin{frame}{SRTF: Shortest-Remaining-Time-First}
  \begin{block}{Definition}
    \key{SRTF} is preemptive SJF. On every arrival, compare the newcomer's
    burst with the \emph{remaining} time of the running process; preempt only
    if the newcomer is strictly shorter.
  \end{block}
  \begin{itemize}
    \item Decisions at every arrival and completion, not only at block/exit.
    \item Equal remaining times do \key{not} preempt: ties go to the earlier
      arrival.
    \item WT now includes time spent back in the ready queue after
      preemption, so $\mathrm{WT}\ne\mathrm{RT}$ in general.
    \item SRTF minimises average WT among preemptive policies when burst
      times are known --- at the cost of more dispatches.
  \end{itemize}
  \muted{Silberschatz et al., Ch.~5 \S5.3.2; Stallings calls it SRT,
  Ch.~9 \S9.2.}
\end{frame}

\begin{frame}{SRTF on $W$: The Numbers}
  \centering
  \footnotesize
  \gantt{0.6}{P1/0/2/gP1,P2/2/4/gP2,P3/4/5/gP3,P2/5/7/gP2,P4/7/11/gP4,%
    P5/11/14/gP5,P1/14/19/gP1}{P1/0,P2/2,P3/4,P4/5,P5/9}
  \\[0.1cm]
  {\footnotesize
  \begin{tabular}{lrrrrrr}
    \metrichead
    \texttt{P1} & 0 & 7 & 19 & 19 & 12 & 0 \\
    \texttt{P2} & 2 & 4 & 7 & 5 & 1 & 0 \\
    \texttt{P3} & 4 & 1 & 5 & 1 & 0 & 0 \\
    \texttt{P4} & 5 & 4 & 11 & 6 & 2 & 2 \\
    \texttt{P5} & 9 & 3 & 14 & 5 & 2 & 2 \\
    \midrule
    \multicolumn{4}{@{}l}{\textbf{Average}} & 7.2 & 3.4 & 0.8 \\
    \bottomrule
  \end{tabular}}
  \vspace{0.1cm}
  \begin{itemize}
    \item At $t=4$, \texttt{P3} (1 ms) beats \texttt{P2}'s remaining 2 ms:
      \texttt{P2} is preempted and resumes at 5.
    \item Average WT falls to 3.4 ms, but \texttt{P1}'s WT rises to 12 ms and
      the schedule needs more switches.
  \end{itemize}
\end{frame}

\begin{frame}{Priority Scheduling}
  \begin{block}{Definition}
    \key{Priority scheduling} gives the CPU to the ready process with the
    highest priority; equal priorities are broken FCFS. It has
    non-preemptive and preemptive forms.
  \end{block}
  \begin{itemize}
    \item \textbf{Internal} priority: a measured quantity --- time limit,
      memory need, I/O-to-CPU ratio.
    \item \textbf{External} priority: set outside the OS --- importance,
      department, payment.
    \item SJF is the special case ``priority $=$ predicted next burst''.
    \item \key{Starvation}: a steady stream of higher-priority arrivals can
      postpone a low-priority process forever.
    \item \key{Ageing}: raise a waiting process's priority gradually. With
      priorities $0$--$63$ and ageing one level per 10 ms, a process entering
      at 63 reaches 0 after $63\times10=630$ ms --- its wait is bounded.
  \end{itemize}
  \muted{Silberschatz et al., Ch.~5 \S5.3, p.~192.}
\end{frame}

\begin{frame}{Priority (Non-preemptive) on $W$}
  \centering
  \footnotesize
  \gantt{0.6}{P1/0/7/gP1,P2/7/11/gP2,P3/11/12/gP3,P5/12/15/gP5,P4/15/19/gP4}%
    {P1/0,P2/2,P3/4,P4/5,P5/9}
  \\[0.1cm]
  {\footnotesize
  \begin{tabular}{lrrrrrr}
    \metrichead
    \texttt{P1} & 0 & 7 & 7 & 7 & 0 & 0 \\
    \texttt{P2} & 2 & 4 & 11 & 9 & 5 & 5 \\
    \texttt{P3} & 4 & 1 & 12 & 8 & 7 & 7 \\
    \texttt{P4} & 5 & 4 & 19 & 14 & 10 & 10 \\
    \texttt{P5} & 9 & 3 & 15 & 6 & 3 & 3 \\
    \midrule
    \multicolumn{4}{@{}l}{\textbf{Average}} & 8.8 & 5.0 & 5.0 \\
    \bottomrule
  \end{tabular}}
  \vspace{0.1cm}
  \begin{itemize}
    \item At $t=7$: \texttt{P2}(1) runs; at 11 \texttt{P3}(2) beats
      \texttt{P5}(2) on earlier arrival; then \texttt{P5}, then \texttt{P4}.
  \end{itemize}
\end{frame}

\begin{frame}{Priority (Preemptive) on $W$}
  \centering
  \footnotesize
  \gantt{0.6}{P1/0/2/gP1,P2/2/6/gP2,P3/6/7/gP3,P1/7/9/gP1,P5/9/12/gP5,%
    P1/12/15/gP1,P4/15/19/gP4}{P1/0,P2/2,P3/4,P4/5,P5/9}
  \\[0.1cm]
  {\footnotesize
  \begin{tabular}{lrrrrrr}
    \metrichead
    \texttt{P1} & 0 & 7 & 15 & 15 & 8 & 0 \\
    \texttt{P2} & 2 & 4 & 6 & 4 & 0 & 0 \\
    \texttt{P3} & 4 & 1 & 7 & 3 & 2 & 2 \\
    \texttt{P4} & 5 & 4 & 19 & 14 & 10 & 10 \\
    \texttt{P5} & 9 & 3 & 12 & 3 & 0 & 0 \\
    \midrule
    \multicolumn{4}{@{}l}{\textbf{Average}} & 7.8 & 4.0 & 2.4 \\
    \bottomrule
  \end{tabular}}
  \vspace{0.1cm}
  \begin{itemize}
    \item At $t=2$, \texttt{P2}(1) preempts \texttt{P1}(3); at $t=9$,
      \texttt{P5}(2) preempts \texttt{P1}(3). \texttt{P4}(4) waits longest.
  \end{itemize}
\end{frame}

\begin{frame}{What Preemption Bought --- and Who Paid}
  \centering
  \renewcommand{\arraystretch}{1.2}
  \begin{tabular}{lcccc}
    \toprule
    \textbf{Policy on $W$} & \textbf{avg WT} & \textbf{WT of P2} &
    \textbf{WT of P5} & \textbf{WT of P4} \\
    \midrule
    FCFS & 5.2 & 5 & 7 & 7 \\
    Priority, non-preemptive & 5.0 & 5 & 3 & 10 \\
    Priority, preemptive & \good{4.0} & \good{0} & \good{0} & \bad{10} \\
    \bottomrule
  \end{tabular}
  \vspace{0.25cm}
  \begin{itemize}
    \item The important process \texttt{P2} and the short \texttt{P5} now wait
      0 ms; \texttt{P4} pays with 10 ms.
    \item Non-preemptive priority could not help \texttt{P2}: \texttt{P1}
      already held the CPU when \texttt{P2} arrived.
    \item Priority is a statement about \emph{importance}, not about length.
      Without ageing, a low-priority process can starve.
  \end{itemize}
\end{frame}

\begin{frame}{Round Robin: FCFS With a Timer}
  \begin{block}{Definition}
    \key{Round Robin} gives each process at most one \key{time quantum} $q$.
    If its burst is unfinished when the quantum expires, it is preempted and
    placed at the tail of the circular ready queue.
  \end{block}
  \begin{itemize}
    \item A burst shorter than $q$ releases the CPU voluntarily.
    \item Timer interrupt enforces the quantum --- this is why Lecture 2's
      timer is a scheduling mechanism.
    \item \key{Tie rule:} an arrival at the instant a quantum expires joins
      the queue before the preempted process.
    \item \key{Guarantee:} with $n$ ready processes, each waits at most
      $(n-1)q$ for its next turn, ignoring switch time.
  \end{itemize}
  \muted{Silberschatz et al., Ch.~5 \S5.3, p.~194; Stallings,
  Ch.~9 \S9.2, p.~410.}
\end{frame}

\begin{frame}{RR with $q=3$ on $W$: The Numbers}
  \centering
  \footnotesize
  \gantt{0.6}{P1/0/3/gP1,P2/3/6/gP2,P1/6/9/gP1,P3/9/10/gP3,P4/10/13/gP4,%
    P2/13/14/gP2,P5/14/17/gP5,P1/17/18/gP1,P4/18/19/gP4}%
    {P1/0,P2/2,P3/4,P4/5,P5/9}
  \\[0.1cm]
  {\footnotesize
  \begin{tabular}{lrrrrrr}
    \metrichead
    \texttt{P1} & 0 & 7 & 18 & 18 & 11 & 0 \\
    \texttt{P2} & 2 & 4 & 14 & 12 & 8 & 1 \\
    \texttt{P3} & 4 & 1 & 10 & 6 & 5 & 5 \\
    \texttt{P4} & 5 & 4 & 19 & 14 & 10 & 5 \\
    \texttt{P5} & 9 & 3 & 17 & 8 & 5 & 5 \\
    \midrule
    \multicolumn{4}{@{}l}{\textbf{Average}} & 11.6 & 7.8 & 3.2 \\
    \bottomrule
  \end{tabular}}
  \vspace{0.1cm}
  \begin{itemize}
    \item $q=3$: P1, P2, P1, P3, P4, P2, P5, P1, P4.
    \item Average RT is 3.2 ms --- far better than FCFS's 5.2 ms --- but
      average WT rises to 7.8 ms. RR trades turnaround for responsiveness.
  \end{itemize}
\end{frame}

\begin{frame}{The Quantum Changes Everything}
  \centering
  \renewcommand{\arraystretch}{1.15}
  \begin{tabular}{lccccc}
    \toprule
    \textbf{RR on $W$} & \textbf{avg WT} & \textbf{avg TAT} &
    \textbf{avg RT} & \textbf{switches} \\
    \midrule
    $q=1$ & 6.80 & 10.60 & \good{1.00} & 17 \\
    $q=2$ & 6.20 & 10.00 & 2.00 & 10 \\
    $q=3$ & \bad{7.80} & \bad{11.60} & 3.20 & 8 \\
    $q=4$ & \good{5.00} & \good{8.80} & 4.00 & 5 \\
    $q=5$ & 5.80 & 9.60 & 4.00 & 5 \\
    $q\ge 9$ & 5.20 & 9.00 & 5.20 & 4 \\
    \bottomrule
  \end{tabular}
  \vspace{0.2cm}
  \begin{itemize}
    \item RT rises monotonically with $q$; average WT does \key{not} ---
      $q=3$ is worse than both $q=2$ and $q=4$. It depends on how quantum
      boundaries line up with bursts.
    \item Once $q\ge$ the longest burst (9 ms), RR \emph{is} FCFS.
    \item Overhead: with switch cost $s$, useful fraction is
      $\frac{q}{q+s}$. For $s=0.5$: $q=1\Rightarrow66.7\%$,
      $q=3\Rightarrow85.7\%$, $q=10\Rightarrow95.2\%$.
  \end{itemize}
\end{frame}

\begin{frame}{HRRN: Highest Response Ratio Next}
  \begin{block}{Definition}
    \key{HRRN} is non-preemptive. At each completion, choose the ready
    process with the highest \key{response ratio}
    \[
      R=\frac{W+\tau}{\tau}=1+\frac{W}{\tau},
    \]
    where $W$ is the time already waited and $\tau$ the required service.
  \end{block}
  \begin{itemize}
    \item A short job gets a high $R$ because $\tau$ is small; a long-waiting
      job gets a high $R$ because $W$ is large.
    \item It behaves like SJF for fresh jobs and like FCFS for old ones ---
      which bounds starvation.
    \item Costs: it must know the service time $\tau$, and it recomputes
      $R$ at every completion.
  \end{itemize}
  \muted{Not in every Silberschatz edition; standard in GATE/UG treatments
  of scheduling.}
\end{frame}

\begin{frame}{HRRN on $W$: The Numbers}
  \centering
  \footnotesize
  \gantt{0.6}{P1/0/7/gP1,P3/7/8/gP3,P2/8/12/gP2,P4/12/16/gP4,P5/16/19/gP5}%
    {P1/0,P2/2,P3/4,P4/5,P5/9}
  \\[0.1cm]
  {\footnotesize
  \begin{tabular}{lrrrrrr}
    \metrichead
    \texttt{P1} & 0 & 7 & 7 & 7 & 0 & 0 \\
    \texttt{P2} & 2 & 4 & 12 & 10 & 6 & 6 \\
    \texttt{P3} & 4 & 1 & 8 & 4 & 3 & 3 \\
    \texttt{P4} & 5 & 4 & 16 & 11 & 7 & 7 \\
    \texttt{P5} & 9 & 3 & 19 & 10 & 7 & 7 \\
    \midrule
    \multicolumn{4}{@{}l}{\textbf{Average}} & 8.4 & 4.6 & 4.6 \\
    \bottomrule
  \end{tabular}}
  \vspace{0.1cm}
  \begin{itemize}
    \item At $t=8$: \texttt{P2} $R=(6+4)/4=2.5$ beats \texttt{P4}
      $R=(3+4)/4=1.75$.
    \item At $t=12$: \texttt{P4} $R=(7+4)/4=2.75$ beats \texttt{P5}
      $R=(3+3)/3=2.0$ --- the long waiter finally wins one.
  \end{itemize}
\end{frame}

\begin{frame}{LJF / LRTF: The Deliberately Bad Mirror Image}
  \begin{block}{Definition}
    \key{Longest-Job-First} picks the largest burst; \key{LRTF} is its
    preemptive form, picking the largest remaining time.
  \end{block}
  \begin{itemize}
    \item If SJF minimises average waiting time, LJF maximises it: the
      largest weight is now placed on the largest burst.
    \item It is useful as a contrast case and in real-time systems where the
      longest job has the earliest deadline --- never as a general-purpose
      policy.
    \item It starves short jobs even more aggressively than SJF starves long
      ones.
  \end{itemize}
  \vspace{0.1cm}
  \begin{exampleblock}{Same three jobs, both extremes}
    All arrive at 0; bursts $(3,5,2)$.
    LJF order $5,3,2$: WT $=0,5,8$, average $4.33$ ms.\\
    SJF order $2,3,5$: WT $=0,2,5$, average $2.33$ ms.
  \end{exampleblock}
\end{frame}

\transitionslide{A single queue forces one policy on everyone. Real systems put processes into classes --- and sometimes let the classes move.}

% =============================================================================
% ACT 3 --- QUEUE-BASED POLICIES
% =============================================================================
\section{Queue-Based Policies}

\begin{frame}{Multilevel Queue (MLQ)}
  \centering
  \begin{tikzpicture}[every node/.style={font=\small}]
    % Queue boxes, priority arrow, per-queue policy labels.
    \node[dag-node, minimum width=6cm, minimum height=0.65cm] (q0) at (0,1.5)
      {system processes};
    \node[dag-node, minimum width=6cm, minimum height=0.65cm] (q1) at (0,0.5)
      {interactive processes};
    \node[dag-node, minimum width=6cm, minimum height=0.65cm] (q2) at (0,-0.5)
      {interactive editing};
    \node[dag-node, minimum width=6cm, minimum height=0.65cm] (q3) at (0,-1.5)
      {batch processes};
    \node[text=primary-gold,anchor=west] at (3.3,1.5) {RR};
    \node[text=primary-gold,anchor=west] at (3.3,0.5) {RR};
    \node[text=primary-gold,anchor=west] at (3.3,-0.5) {RR};
    \node[text=primary-gold,anchor=west] at (3.3,-1.5) {FCFS};
    \draw[observed-edge] (-4.3,-1.85) -- (-4.3,1.85);
    \node[align=center,font=\footnotesize] at (-5.5,0) {higher\\priority};
  \end{tikzpicture}
  \vspace{0.05cm}
  \begin{itemize}
    \item Each process is \key{permanently} assigned to one queue by type,
      memory need, or priority.
    \item Each queue runs its own algorithm; a policy \emph{between} queues
      is still required:
      \begin{itemize}
        \item fixed priority, preemptive --- lower queues can starve;
        \item time slicing, e.g.\ 80\% foreground / 20\% background --- no
          starvation, but no strict priority either.
      \end{itemize}
  \end{itemize}
  \muted{After Silberschatz et al., Ch.~5 \S5.3, p.~196 (Fig.~5.6).}
\end{frame}

\begin{frame}{MLQ Worked Example: Two Queues}
  \footnotesize
  \begin{columns}[T]
    \begin{column}{0.28\textwidth}
      \renewcommand{\arraystretch}{1.15}
      \begin{tabular}{lccc}
        \toprule
        & \textbf{AT} & $\tau$ & \textbf{queue} \\
        \midrule
        \texttt{A} & 0 & 4 & fg \\
        \texttt{B} & 1 & 5 & bg \\
        \texttt{C} & 2 & 2 & fg \\
        \texttt{D} & 3 & 3 & bg \\
        \bottomrule
      \end{tabular}
      \vspace{0.2cm}

      Foreground: RR, $q=2$.\\
      Background: FCFS.\\
      Between queues: fixed priority, preemptive.
    \end{column}
    \begin{column}{0.68\textwidth}
      \centering
      \gantt{0.5}{A/0/2/gFg,C/2/4/gFg,A/4/6/gFg,B/6/11/gBg,D/11/14/gBg}{A/0,B/1,C/2,D/3}
    \end{column}
  \end{columns}
  \vspace{0.15cm}
  \begin{itemize}
    \item $t=2$: \texttt{C} arrives in the foreground just as \texttt{A}'s
      quantum expires, so \texttt{C} runs next; \texttt{A} resumes at 4.
    \item $t=6$: foreground empty, so the background starts. \texttt{B}
      (arrived at 1) has waited 5 ms; \texttt{D} waits 8 ms.
    \item WT: A $2$, B $5$, C $0$, D $8$; average $15/4=3.75$ ms.
      Foreground never waits for background.
  \end{itemize}
\end{frame}

\begin{frame}{Multilevel Feedback Queue (MLFQ)}
  \centering
  \begin{tikzpicture}[every node/.style={font=\small}]
    \node[dag-node, minimum width=5.4cm, minimum height=0.7cm] (q0) at (0,1.5)
      {$Q_0$: RR, $q=2$ ms};
    \node[dag-node, minimum width=5.4cm, minimum height=0.7cm] (q1) at (0,0)
      {$Q_1$: RR, $q=4$ ms};
    \node[dag-node, minimum width=5.4cm, minimum height=0.7cm] (q2) at (0,-1.5)
      {$Q_2$: FCFS};
    \draw[observed-edge] (q0.east) -- (3.8,1.5) -- (3.8,0) -- (q1.east);
    \draw[observed-edge] (q1.east) -- (3.8,0) -- (3.8,-1.5) -- (q2.east);
    \node[anchor=west,font=\scriptsize,text=primary-gold] at (3.95,0.75)
      {used full quantum};
    \node[anchor=west,font=\scriptsize,text=primary-gold] at (3.95,-0.75)
      {used full quantum};
    \draw[observed-edge] (-4.9,1.5) -- (q0.west);
    \node[anchor=east,font=\scriptsize] at (-5.0,1.5) {new arrival};
  \end{tikzpicture}
  \vspace{0.1cm}
  \begin{itemize}
    \item Processes \key{move}: a full quantum demotes one level; an
      I/O-bound process that finishes early stays high.
    \item Promotion (ageing) prevents starvation in the bottom queue.
    \item With the right parameters MLFQ behaves like RR, like FCFS, or like
      an \emph{approximation of SJF} --- without ever being told $\tau$.
  \end{itemize}
  \muted{After Silberschatz et al., Ch.~5 \S5.3, p.~197.}
\end{frame}

\begin{frame}{MLFQ Worked Example: A Late Short Job Wins}
  \centering
  \footnotesize
  Queues: $Q_0$ RR $q=2$, $Q_1$ RR $q=4$, $Q_2$ FCFS; arrivals enter $Q_0$;
  a full quantum demotes.\\[0.1cm]
  \gantt{0.6}{P1/0/2/gQ0,P2/2/4/gQ0,P3/4/6/gQ0,P1/6/8/gQ1,P4/8/10/gQ0,%
    P2/10/11/gQ1,P3/11/15/gQ1,P1/15/19/gQ1,P1/19/20/gQ2}{P1/0,P2/1,P3/2,P4/8}
  \\[0.1cm]
  {\scriptsize
  \begin{tabular}{lrrrrrr}
    \metrichead
    \texttt{P1} & 0 & 9 & 20 & 20 & 11 & 0 \\
    \texttt{P2} & 1 & 3 & 11 & 10 & 7 & 1 \\
    \texttt{P3} & 2 & 6 & 15 & 13 & 7 & 2 \\
    \texttt{P4} & 8 & 2 & 10 & 2 & 0 & 0 \\
    \midrule
    \multicolumn{4}{@{}l}{\textbf{Average}} & 11.25 & 6.25 & 0.75 \\
    \bottomrule
  \end{tabular}}
  \vspace{0.1cm}
  \begin{itemize}
    \item 0--6: every process gets one 2 ms slice in $Q_0$; none finishes,
      so all three are demoted.
    \item 8--10: \texttt{P4} (2 ms, arrives at 8) runs in $Q_0$ and finishes
      at 10 --- \good{a 2 ms turnaround for a late arrival}.
    \item 10--15: $Q_1$ drains \texttt{P2} (1 ms left) and \texttt{P3}
      (4 ms). \texttt{P1} then runs 15--19 in $Q_1$, is demoted, and finishes
      its last 1 ms in $Q_2$ at 20.
  \end{itemize}
\end{frame}

\begin{frame}{MLFQ: Five Knobs and One Way to Game It}
  \begin{block}{The five parameters}
    \begin{enumerate}
      \item number of queues;
      \item scheduling algorithm per queue;
      \item demotion rule;
      \item promotion (ageing) rule;
      \item entry queue for a new process.
    \end{enumerate}
  \end{block}
  \begin{itemize}
    \item \key{Gaming:} a CPU-bound process can issue a tiny I/O request
      just before each quantum expires, so it never ``uses a full quantum''
      and never gets demoted.
    \item \key{Defence:} charge CPU time cumulatively at each level, and
      demote when the allotment is exhausted however it was split.
    \item Periodic promotion protects the processes the gamer would
      otherwise starve.
  \end{itemize}
  \muted{Silberschatz et al., Ch.~5 \S5.3, p.~197.}
\end{frame}

% =============================================================================
% ACT 4 --- COMPARISON AND CHOICE
% =============================================================================
\section{Comparison and Choice}

\transitionslide{Seven policies, one workload. If one algorithm won every column, scheduling would be an easy subject.}

\begin{frame}{All Seven Policies on $W$}
  \centering
  \footnotesize
  \renewcommand{\arraystretch}{1.22}
  \begin{tabular}{lcccc}
    \toprule
    \textbf{Policy} & \textbf{avg WT} & \textbf{avg TAT} &
    \textbf{avg RT} & \textbf{switches} \\
    \midrule
    FCFS & 5.2 & 9.0 & 5.2 & 4 \\
    SJF (non-preemptive) & 4.4 & 8.2 & 4.4 & 4 \\
    SRTF & \good{3.4} & \good{7.2} & \good{0.8} & 6 \\
    Priority (non-preemptive) & 5.0 & 8.8 & 5.0 & 4 \\
    Priority (preemptive) & 4.0 & 7.8 & 2.4 & 6 \\
    RR, $q=3$ & \bad{7.8} & \bad{11.6} & 3.2 & 8 \\
    HRRN & 4.6 & 8.4 & 4.6 & 4 \\
    \bottomrule
  \end{tabular}
  \vspace{0.2cm}
  \begin{itemize}
    \item SRTF wins all three averages here, but it needs burst prediction,
      causes the most starvation, and pays for six dispatches.
    \item RR is worst on WT and TAT and best on nothing here --- except the
      guarantee that no one waits more than $(n-1)q$.
    \item ``Best'' is always relative to the criterion the system values.
  \end{itemize}
  \muted{All times in ms. Switch count = changes of the running process.}
\end{frame}

\begin{frame}{How to Choose: A One-Minute Decision Guide}
  \centering
  \footnotesize
  \renewcommand{\arraystretch}{1.22}
  \begin{tabular}{p{5.3cm}p{6.8cm}}
    \toprule
    \textbf{If the system is \ldots} & \textbf{Reach for \ldots} \\
    \midrule
    batch, no interaction, simple & FCFS (or SJF if bursts are known) \\
    interactive / time-sharing & RR, or MLFQ for mixed behaviour \\
    priority-driven (kernel threads, real-time-ish) & Priority + ageing \\
    fixed classes (system, interactive, batch) & MLQ with a between-queue rule \\
    SJF-like, but bursts unknown & MLFQ \\
    service times known, long jobs must not starve & HRRN \\
    \bottomrule
  \end{tabular}
  \vspace{0.3cm}
  \begin{block}{Exam translation table}
    ``preemptive SJF'' $=$ SRTF \quad\textbullet\quad
    ``time slice'' $=$ RR quantum \quad\textbullet\quad
    ``quantum $\to\infty$'' $=$ FCFS \quad\textbullet\quad
    ``feedback'' $=$ processes migrate between queues
  \end{block}
\end{frame}

\begin{frame}{Common Exam Traps}
  \footnotesize
  \begin{columns}[T]
    \begin{column}{0.48\textwidth}
      \begin{block}{Wrong}
        \begin{itemize}
          \item ``Waiting time is CT $-$ AT.''
          \item ``Response time is when the process finishes.''
          \item ``SJF is always optimal.''
          \item ``A bigger quantum always lowers WT.''
          \item ``SRTF preempts on equal remaining time.''
          \item ``MLQ moves a process between queues.''
          \item ``Idle CPU time counts as someone's waiting time.''
        \end{itemize}
      \end{block}
    \end{column}
    \begin{column}{0.48\textwidth}
      \begin{block}{Right}
        \begin{itemize}
          \item That is TAT; WT $=$ TAT $-$ BT.
          \item RT stops at the first dispatch.
          \item Optimal only for a simultaneous batch, and only for WT.
          \item WT is not monotone in $q$ (6.2, 7.8, 5.0 on $W$).
          \item Only a strictly shorter newcomer preempts.
          \item Only MLFQ migrates; MLQ assignment is permanent.
          \item No process is ready, so nobody waited.
        \end{itemize}
      \end{block}
    \end{column}
  \end{columns}
\end{frame}

% =============================================================================
% ACT 5 --- GATE-PATTERN HOMEWORK
% =============================================================================
\section{GATE-Pattern Homework}

\transitionslide{Homework: six numerical problems and one design question, in GATE style. Attempt first; solutions follow.}

\begin{frame}{Homework Set: Instructions}
  \footnotesize
  \begin{block}{Before you turn the page}
    \begin{itemize}
      \item Time: 45 minutes. Calculator not needed; all numbers are small.
      \item Use the tie rules from this lecture: earlier arrival, then lower
        index; SRTF preempts only if strictly shorter; RR arrivals at a
        quantum boundary join before the preempted process.
      \item For every question: draw the Gantt chart \emph{first}, then fill
        CT, TAT, WT, RT, then average.
    \end{itemize}
  \end{block}
  \begin{itemize}
    \item Q1: three policies on one batch.
    \item Q2: SRTF trace.
    \item Q3: RR with $q=3$.
    \item Q4: five true/false statements.
    \item Q5: HRRN versus SJF.
    \item Q6: two-queue MLQ.
    \item Q7: design choices.
  \end{itemize}
\end{frame}

\begin{frame}{Q1--Q2}
  \begin{block}{Q1 \normalfont (metrics, 8 marks)}
    Three processes arrive at time 0 in the order \texttt{P1}, \texttt{P2},
    \texttt{P3} with bursts $8,4,2$ ms. Compute average WT and average TAT
    under (a) FCFS, (b) non-preemptive SJF, (c) RR with $q=2$. Which policy
    minimises average TAT?
  \end{block}
  \begin{block}{Q2 \normalfont (SRTF, 8 marks)}
    \texttt{P1}$(0,8)$, \texttt{P2}$(1,4)$, \texttt{P3}$(2,9)$,
    \texttt{P4}$(3,5)$ as (AT, $\tau$) in ms. Draw the SRTF Gantt chart,
    find the average waiting time, and state \texttt{P1}'s completion time.
  \end{block}
\end{frame}

\begin{frame}{Q3--Q4}
  \begin{block}{Q3 \normalfont (RR, 8 marks)}
    Five processes arrive together with bursts $6,4,2,8,3$ ms; RR uses
    $q=3$. Give the completion order, the average WT, and the number of
    time slices used.
  \end{block}
  \begin{block}{Q4 \normalfont (true/false, 10 marks)}
    State whether each is TRUE or FALSE and justify in one line:
    \begin{enumerate}
      \item[(i)] If all processes arrive at $t=0$, non-preemptive SJF
        minimises average WT.
      \item[(ii)] Preemptive priority scheduling can starve a low-priority
        process.
      \item[(iii)] RR with $q$ larger than every burst is FCFS.
      \item[(iv)] In SRTF, a newcomer with remaining time equal to the
        running process preempts it.
      \item[(v)] Increasing the RR quantum always decreases average WT.
    \end{enumerate}
  \end{block}
\end{frame}

\begin{frame}{Q5--Q7}
  \begin{block}{Q5 \normalfont (HRRN, 8 marks)}
    \texttt{P1}$(0,10)$, \texttt{P2}$(1,6)$, \texttt{P3}$(10,5)$ as
    (AT, $\tau$). (a) Schedule by HRRN: show the response ratios at $t=10$.
    (b) Compare average WT with non-preemptive SJF.
  \end{block}
  \begin{block}{Q6 \normalfont (MLQ, 8 marks)}
    Two queues: foreground RR $q=2$, background FCFS, fixed priority
    preemptive. \texttt{A}$(0,5)$ and \texttt{C}$(2,3)$ are foreground;
    \texttt{B}$(1,4)$ and \texttt{D}$(4,2)$ are background. Draw the Gantt
    chart, find average WT, and explain why \texttt{B} waits despite
    arriving before \texttt{C}.
  \end{block}
  \begin{block}{Q7 \normalfont (design, 5 marks)}
    Name the best policy for: (a) a nightly payroll batch; (b) an
    interactive shell; (c) a mix of system, interactive, and batch work;
    (d) SJF-like behaviour when bursts are unknown.
  \end{block}
\end{frame}

\begin{frame}{Solution 1: Three Policies on One Batch}
  \centering
  \footnotesize
  \gantt{0.5}{P1/0/8/gP1,P2/8/12/gP2,P3/12/14/gP3}{}
  \quad
  \gantt{0.5}{P3/0/2/gP3,P2/2/6/gP2,P1/6/14/gP1}{}
  \\[0.05cm]
  \gantt{0.6}{P1/0/2/gP1,P2/2/4/gP2,P3/4/6/gP3,P1/6/8/gP1,P2/8/10/gP2,%
    P1/10/14/gP1}{}
  \\[0.15cm]
  {\footnotesize
  \begin{tabular}{lcccl}
    \toprule
    \textbf{Policy} & \textbf{avg WT} & \textbf{avg TAT} &
    \textbf{avg RT} & \textbf{Why} \\
    \midrule
    FCFS & 6.67 & 11.33 & 6.67 & order 1,2,3 \\
    SJF & \good{2.67} & \good{7.33} & 2.67 & order 3,2,1 \\
    RR $q=2$ & 5.33 & 10.00 & \good{2.00} & six slices \\
    \bottomrule
  \end{tabular}}
  \vspace{0.1cm}
  \begin{itemize}
    \item SJF minimises average TAT. RR gives the best average RT, but its
      WT sits between FCFS and SJF.
  \end{itemize}
\end{frame}

\begin{frame}{Solution 2: SRTF Trace}
  \centering
  \footnotesize
  \gantt{0.48}{P1/0/1/gP1,P2/1/5/gP2,P4/5/10/gP4,P1/10/17/gP1,P3/17/26/gP3}%
    {P1/0,P2/1,P3/2,P4/3}
  \\[0.1cm]
  {\footnotesize
  \begin{tabular}{lrrrrrr}
    \metrichead
    \texttt{P1} & 0 & 8 & 17 & 17 & 9 & 0 \\
    \texttt{P2} & 1 & 4 & 5 & 4 & 0 & 0 \\
    \texttt{P3} & 2 & 9 & 26 & 24 & 15 & 15 \\
    \texttt{P4} & 3 & 5 & 10 & 7 & 2 & 2 \\
    \midrule
    \multicolumn{4}{@{}l}{\textbf{Average}} & 13.0 & \good{6.5} & 4.25 \\
    \bottomrule
  \end{tabular}}
  \vspace{0.1cm}
  \begin{itemize}
    \item At $t=1$, \texttt{P2}(4) beats \texttt{P1}'s remaining 7; at $t=3$,
      \texttt{P4}(5) does \emph{not} beat \texttt{P2}'s remaining 2.
    \item \texttt{P1} completes at \good{17 ms}; average WT $=26/4=6.5$ ms.
  \end{itemize}
\end{frame}

\begin{frame}{Solution 3: RR with $q=3$}
  \centering
  \footnotesize
  \gantt{0.5}{P1/0/3/gP1,P2/3/6/gP2,P3/6/8/gP3,P4/8/11/gP4,P5/11/14/gP5,%
    P1/14/17/gP1,P2/17/18/gP2,P4/18/23/gP4}{}
  \\[0.1cm]
  {\footnotesize
  \begin{tabular}{lrrrrrr}
    \metrichead
    \texttt{P1} & 0 & 6 & 17 & 17 & 11 & 0 \\
    \texttt{P2} & 0 & 4 & 18 & 18 & 14 & 3 \\
    \texttt{P3} & 0 & 2 & 8 & 8 & 6 & 6 \\
    \texttt{P4} & 0 & 8 & 23 & 23 & 15 & 8 \\
    \texttt{P5} & 0 & 3 & 14 & 14 & 11 & 11 \\
    \midrule
    \multicolumn{4}{@{}l}{\textbf{Average}} & 16.0 & \bad{11.4} & 5.6 \\
    \bottomrule
  \end{tabular}}
  \vspace{0.1cm}
  \begin{itemize}
    \item Completion order: \texttt{P3}(8), \texttt{P5}(14), \texttt{P1}(17),
      \texttt{P2}(18), \texttt{P4}(23).
    \item Slices: \texttt{P1},\texttt{P2},\texttt{P3},\texttt{P4},
      \texttt{P5},\texttt{P1},\texttt{P2},\texttt{P4} --- \good{8 slices}.
  \end{itemize}
\end{frame}

\begin{frame}{Solution 4: Five Statements}
  \centering
  \footnotesize
  \renewcommand{\arraystretch}{1.3}
  \begin{tabular}{clp{8.2cm}}
    \toprule
    & \textbf{Answer} & \textbf{Reason} \\
    \midrule
    (i) & \good{TRUE} & All ready at 0: the exchange argument applies. \\
    (ii) & \good{TRUE} & A stream of higher-priority arrivals postpones it
      indefinitely. \\
    (iii) & \good{TRUE} & No process is ever preempted; the queue is FIFO. \\
    (iv) & \bad{FALSE} & SRTF preempts only on a \emph{strictly} shorter
      remaining time. \\
    (v) & \bad{FALSE} & WT is not monotone in $q$; on $W$, $q=3$ is worse
      than $q=2$ and $q=4$. \\
    \bottomrule
  \end{tabular}
  \vspace{0.3cm}
  \begin{itemize}
    \item The two false statements are the two most common GATE traps:
      non-strict preemption and a monotonicity assumption that does not hold.
  \end{itemize}
\end{frame}

\begin{frame}{Solution 5: HRRN versus SJF}
  \centering
  \footnotesize
  \textbf{HRRN:} \gantt{0.5}{P1/0/10/gP1,P2/10/16/gP2,P3/16/21/gP3}%
    {P1/0,P2/1,P3/10}
  \\[0.1cm]
  \textbf{SJF:} \gantt{0.5}{P1/0/10/gP1,P3/10/15/gP3,P2/15/21/gP2}%
    {P1/0,P2/1,P3/10}
  \\[0.1cm]
  {\footnotesize
  \begin{tabular}{lccc}
    \toprule
    \textbf{Policy} & \textbf{avg WT} & \textbf{WT of P2} & \textbf{WT of P3} \\
    \midrule
    HRRN & 5.00 & \good{9} & 6 \\
    SJF & \good{4.67} & \bad{14} & \good{0} \\
    \bottomrule
  \end{tabular}}
  \vspace{0.1cm}
  \begin{itemize}
    \item At $t=10$: $R_2=(9+6)/6=2.5$, $R_3=(0+5)/5=1.0$ --- HRRN picks
      \texttt{P2}.
    \item SJF wins average WT by 0.33 ms, but \texttt{P2} waits 14 ms
      instead of 9 ms. HRRN converts a small average loss into a bounded
      wait for the long waiter.
  \end{itemize}
\end{frame}

\begin{frame}{Solution 6: Two-Queue MLQ}
  \centering
  \footnotesize
  \gantt{0.6}{A/0/2/gFg,C/2/4/gFg,A/4/6/gFg,C/6/7/gFg,A/7/8/gFg,%
    B/8/12/gBg,D/12/14/gBg}{A/0,B/1,C/2,D/4}
  \\[0.1cm]
  {\footnotesize
  \begin{tabular}{lrrrrrr}
    \metrichead
    \texttt{A} & 0 & 5 & 8 & 8 & 3 & 0 \\
    \texttt{B} & 1 & 4 & 12 & 11 & 7 & 7 \\
    \texttt{C} & 2 & 3 & 7 & 5 & 2 & 0 \\
    \texttt{D} & 4 & 2 & 14 & 10 & 8 & 8 \\
    \midrule
    \multicolumn{4}{@{}l}{\textbf{Average}} & 8.5 & \good{5.0} & 3.75 \\
    \bottomrule
  \end{tabular}}
  \vspace{0.1cm}
  \begin{itemize}
    \item Foreground slices: \texttt{A}(0--2), \texttt{C}(2--4),
      \texttt{A}(4--6), \texttt{C}(6--7), \texttt{A}(7--8). Then background
      FCFS: \texttt{B}, \texttt{D}.
    \item \texttt{B} arrives at 1 but is background; fixed priority says
      foreground first, so it waits until 8. A steady foreground stream
      would starve it --- the reason MLQ needs a time-slicing rule or ageing.
  \end{itemize}
\end{frame}

\begin{frame}{Solution 7: Design Choices}
  \centering
  \footnotesize
  \renewcommand{\arraystretch}{1.35}
  \begin{tabular}{p{5.6cm}p{6.2cm}}
    \toprule
    \textbf{Workload} & \textbf{Best fit and why} \\
    \midrule
    (a) nightly payroll batch & FCFS: simple, no user waiting, job order
      irrelevant. SJF if runtimes are known. \\
    (b) interactive shell & RR (small $q$): bounds response time; FCFS would
      freeze the user behind a long job. \\
    (c) system + interactive + batch & MLQ: fixed classes, each with its own
      policy; add time slicing between queues to avoid starvation. \\
    (d) SJF-like, bursts unknown & MLFQ: infers short jobs by demoting
      full-quantum users; no burst prediction needed. \\
    \bottomrule
  \end{tabular}
  \vspace{0.25cm}
  \begin{itemize}
    \item Full marks require the \emph{reason}, not just the name: what the
      workload values (response, throughput, fairness) and what information
      the scheduler actually has.
  \end{itemize}
\end{frame}

% =============================================================================
% ACT 6 --- WRAP-UP
% =============================================================================
\section{Wrap-Up}

\begin{frame}{Takeaways}
  \begin{itemize}
    \item Measure before comparing: TAT $=$ CT $-$ AT, WT $=$ TAT $-$ BT,
      RT $=$ ST $-$ AT; utilization and throughput are system-wide.
    \item \key{FCFS}: simple, starvation-free, convoy-prone.
    \item \key{SJF}: optimal average WT for a ready batch, but needs a burst
      prediction; \key{SRTF} is its preemptive form.
    \item \key{Priority}: encodes importance; preemption helps important
      arrivals; ageing bounds starvation.
    \item \key{RR}: bounds response at $(n-1)q$; the quantum trades RT
      against WT, TAT, and switch overhead.
    \item \key{HRRN}: SJF-like, but the response ratio protects long waiters.
    \item \key{MLQ}: fixed classes, no migration; \key{MLFQ}: learned
      classes, migration, and the closest practical approximation to SJF.
    \item Always draw the Gantt chart first, then compute per process, then
      average. Never average an average of something else.
  \end{itemize}
\end{frame}

\begin{frame}{Further Reading (Page-Cited)}
  \footnotesize
  \begin{itemize}
    \item A.~Silberschatz, P.~B.~Galvin, G.~Gagne, \emph{Operating System
      Concepts}, 8e --- Ch.~5 \S5.1 (basic concepts, p.~183), \S5.2
      (criteria, p.~187), \S5.3.1 (FCFS, p.~188), \S5.3.2 (SJF/SRTF,
      p.~189), \S5.3 (priority p.~192, RR p.~194, MLQ p.~196, MLFQ p.~197),
      \S5.5 (multiprocessor, p.~200).
    \item W.~Stallings, \emph{Operating Systems: Internals and Design
      Principles}, 6e --- Ch.~9 \S9.1 (types of scheduling, p.~406), \S9.2
      (FCFS, SPN, SRT, RR, priority, HRRN, p.~410 ff.).
    \item D.~P.~Bovet \& M.~Cesati, \emph{Understanding the Linux Kernel},
      3e --- Ch.~7 (process scheduling, pp.~258--293).
  \end{itemize}
  \vspace{0.2cm}
  \muted{Page-verified via each book's \texttt{index.md} under
  \texttt{master\_supporting\_docs/CS502/supporting\_books/}.}
\end{frame}

\begin{frame}[allowframebreaks]{References}
  \footnotesize
  \bibliographystyle{plain}
  \nocite{SilberschatzOSC_operating_system_concepts,
    StallingsOS_internals_design_principles,
    BovetCesati_understanding_linux_kernel}
  \bibliography{../../Bibliography_base}
\end{frame}

\end{document}
